\documentclass[9pt,a4paper]{extarticle}
\usepackage[french]{babel}
\usepackage[margin=0.8cm]{geometry}
\usepackage{multicol}
\usepackage{enumitem}
\usepackage{titlesec}
\usepackage{xcolor}
\usepackage{fontspec}

\setmainfont{Segoe UI}

\setlength{\columnsep}{0.4cm}
\setlength{\parindent}{0pt}
\setlength{\parskip}{2pt}

\titleformat{\section}{\normalfont\bfseries\color{blue}}{\thesection}{0.5em}{}
\titleformat{\subsection}{\normalfont\bfseries\small\color{teal}}{\thesubsection}{0.5em}{}
\titlespacing*{\section}{0pt}{5pt}{3pt}
\titlespacing*{\subsection}{0pt}{4pt}{2pt}

\setlist[itemize]{noitemsep, topsep=0pt, leftmargin=*}

\pagestyle{empty}

\begin{document}

\begin{center}
\textbf{\large FICHE DE RÉVISION - GESTION D'INTRUSION}
\end{center}

\begin{multicols}{2}

\section*{1. ORIGINE CYBERSÉCURITÉ}

\subsection*{Premières Menaces (1970-1980)}
\textbf{1971 - Creeper:} Premier programme auto-réplicatif sur ARPANET. Bob Thomas (BBN). Message: "I'm the creeper, catch me if you can!". Non malveillant, expérimentation.

\textbf{1988 - Morris Worm:} Premier ver grande échelle. Robert Tappan Morris. 6000 machines (~10\% Internet). Exploitait: sendmail, finger, mots de passe UNIX. Premier incident majeur mondial.

\textbf{1987 - Premiers Antivirus:} Détection virus, signatures statiques, protection locale.

\subsection*{Facteurs Déclencheurs}
Interconnexion systèmes (ARPANET vers Internet académique). Standardisation TCP/IP (1983). Militarisation réseaux (Défense US, résilience). Triade CIA: Confidentialité, Intégrité, Disponibilité. Besoin confidentialité: données militaires, universitaires, bancaires.

\subsection*{Premières Réponses (1988)}
\textbf{CERT:} Computer Emergency Response Team, suite Morris Worm.

\textbf{Pare-feu:} Filtrage IP/port, contrôle flux, architectures périmétriques.

\textbf{Authentification:} RSA 1977, politiques mots de passe, identification centralisée.

\subsection*{Évolution Chronologique}
\textbf{1990s - Périmétrique:} Château fort, pare-feu, DMZ.

\textbf{2000s - Défense profondeur:} IDS/IPS, antivirus, segmentation, anomalies.

\textbf{2010-2020 - Moderne:} Zero Trust, Security by Design, Threat Intelligence, DevSecOps.

\textbf{Actuel:} IA/ML détection, cloud-native, XDR (Extended Detection and Response).

\section*{2. CONCEPTS FONDAMENTAUX}

\subsection*{Détection d'Intrusion - Définition}
Processus identification activité malveillante/anormale dans système, réseau, application.

\textbf{Objectifs:} Détecter attaques en cours, identifier comportements suspects, repérer violations politique, réduire MTTD (Mean Time To Detect).

\textbf{Pourquoi essentiel?} Pare-feu insuffisants, attaques contournent périmètre, APT invisibles semaines, détection = 2ème ligne défense.

\textbf{Positionnement Kill Chain:} Reconnaissance, Exploitation, C2, Mouvement latéral, Exfiltration.

\textbf{Objectif stratégique:} Réactive vers Proactive (détection précoce).

\subsection*{3 Types de Détection}
\textbf{1. Par signature:} Règles connues, rapide, efficace attaques connues, inefficace zero-day. Ex: règle Snort scan Nmap.

\textbf{2. Comportementale (Anomalies):} Modèle comportement normal, écarts statistiques, adaptée APT.

\textbf{3. IA/Machine Learning:} Classification (RF, KNN, CNN, LSTM), patterns complexes, sensible attaques adversariales.

\subsection*{Systèmes de Détection}
\textbf{NIDS:} Analyse trafic réseau.

\textbf{HIDS:} Surveillance hôte.

\textbf{IPS:} Prévention inline (bloque).

\textbf{SIEM:} Corrélation \& vision globale.

\textbf{Cycle:} Collecte - Analyse - Alerte - Investigation - Réponse

\subsection*{Réponse à Incident - Définition}
Ensemble processus organisationnels, techniques, juridiques pour: contenir incident, limiter impact, restaurer systèmes, préserver preuves, éviter récidive.

\textbf{Pourquoi critique?} Détection différent de protection totale, compromis peut être actif, temps réaction = impact.

\textbf{Objectif:} Réduire MTTR (Mean Time To Respond).

\textbf{Cadres:} NIST SP 800-61, DGSSI Maroc, ISO/IEC 27035.

\textbf{Cycle cyber:} Détection - Réponse - Éradication - Restauration - REX

\subsection*{Cycle NIST (6 phases)}
\textbf{1. Préparation:} Politique IR, Playbooks, SOC, CERT interne.

\textbf{2. Détection \& Analyse:} Qualification alerte, Triage L1/L2, forensic, MITRE ATT\&CK.

\textbf{3. Confinement:} Isolation machine, Blocage IP/C2, Segmentation, Suspension comptes.

\textbf{4. Éradication:} Suppression malware, Patch, Rotation MDP, Nettoyage persistence.

\textbf{5. Récupération:} Restauration backups, Vérification intégrité, Surveillance renforcée.

\textbf{6. REX:} Analyse causes, MAJ règles IDS, Amélioration playbooks, Sensibilisation.

\textbf{Rôles SOC:} L1 (triage), L2/L3 (investigation), Incident Manager, Forensic, RSSI.

\subsection*{Gestion Vulnérabilités - Définition}
Processus continu: Identifier, Évaluer, Prioriser, Corriger, Surveiller faiblesses techniques.

\textbf{Vulnérabilité:} Faiblesse exploitable par menace (CIA).

\textbf{Pourquoi stratégique?} 80-90\% intrusions = vulnérabilités connues. Attaques réussies = patch manquant. Mesure préventive majeure.

\textbf{Vulnérabilité différent de Menace différent de Exploit:} Vulnérabilité = faiblesse technique. Exploit = code/méthode exploitation. Menace = acteur/événement malveillant.

\subsection*{Vulnerability Management (5 étapes)}
\textbf{1. Découverte actifs:} Inventaire IT/OT, Shadow IT, Cartographie.

\textbf{2. Scan vulnérabilités:} Nessus, OpenVAS, Qualys, CVE, authentifié/non-auth.

\textbf{3. Évaluation \& Priorisation:} Score CVSS, criticité métier, exposition Internet, exploit public - Risk-Based.

\textbf{4. Remédiation:} Patch management, Reconfiguration, Désactivation, Segmentation.

\textbf{5. Vérification \& Reporting:} Re-scan, KPI (MTTR vuln, \% critiques, conformité patch).

\textbf{Normes:} ISO 27001-A.12.6.1, NIST CSF, DGSSI Maroc.

\subsection*{Sécurité Défensive vs Offensive}

\textbf{Blue Team (Défensive):} Protéger, détecter, répondre, maintenir résilience. Outils: IDS/IPS, SIEM, SOC, EDR. Objectif: Réduire impact attaque. KPI: MTTD/MTTR.

\textbf{Red Team (Offensive):} Simuler attaquant réel, identifier failles, tester défenses. Outils: Metasploit, Nmap, Cobalt Strike. Objectif: Découvrir faiblesses avant attaquant. KPI: Taux compromission.

\textbf{Purple Team:} Collaboration Blue + Red + Threat Intelligence. Objectif: Améliorer posture sécurité globale.

\textbf{Cycle amélioration:} Red attaque - Blue détecte - Analyse écarts - MAJ règles - Renforcement.

\subsection*{Sûreté de Fonctionnement (FMCD)}
Fiabilité (continuité), Maintenabilité (réparation), Confidentialité (non-divulgation), Disponibilité (accessibilité), Intégrité (non-altération).

\textbf{Convergence:} Cyberattaques menacent disponibilité (DoS), intégrité compromise = défaillances. Approche holistique nécessaire.

\subsection*{Statistiques 2023-2024}
93\% entreprises subissent cyberattaque | 4,45 M\$ coût moyen violation | 277 jours détection intrusion | 43\% attaques ciblent PME

\textit{"Il ne s'agit pas de savoir SI vous serez attaqué, mais QUAND."}

\section*{3. CLASSIFICATION ATTAQUES}

\subsection*{Cyber Kill Chain (Lockheed Martin 2011)}
\textbf{7 Phases:}

1. \textbf{Reconnaissance:} Scanning, harvesting infos

2. \textbf{Weaponization:} Création malware/exploit

3. \textbf{Delivery:} Livraison (email, USB, web)

4. \textbf{Exploitation:} Déclenchement code

5. \textbf{Installation:} Installation malware

6. \textbf{Command \& Control (C2):} Canal commande

7. \textbf{Actions on Objectives:} Exfiltration, sabotage

\textbf{Catégories:} 1-3 = Intrusion | 4-5 = Incident | 6-7 = Compromission

\subsection*{Intrusion (Tentative)}
Tentative accès non autorisé. Peut être bloquée AVANT succès. Détectable par systèmes sécurité. N'implique pas nécessairement impact.

\textbf{Ex:} Scan ports, Brute force SSH, Injection SQL bloquée.

\subsection*{Incident (Impact CIA)}
Événement compromettant CIA.

\textbf{Confidentialité:} Vol données clients.

\textbf{Intégrité:} Défacement site web.

\textbf{Disponibilité:} Attaque DDoS.

\subsection*{Compromission (Contrôle Effectif)}
Prise contrôle EFFECTIVE système. Persistance attaquant. Accès privilégié obtenu. Capacité action sur système.

\textbf{Ex:} Backdoor actif, Compte admin créé, C2 établi.

\subsection*{Progression Gravité}
\textbf{Intrusion} (tentative, bloquable) puis \textbf{Incident} (succès, impact) puis \textbf{Compromission} (contrôle total, persistance)

\subsection*{APT - Advanced Persistent Threat}
Attaque sophistiquée, ciblée, prolongée par acteurs étatiques.

\textbf{A}dvanced: Techniques sophistiquées, zero-day.

\textbf{P}ersistent: Présence longue durée (mois/années).

\textbf{T}hreat: Objectifs stratégiques (espionnage, sabotage).

\textbf{Exemples:} APT28, APT29 (Russie), Lazarus Group (Corée du Nord).

\section*{POINTS CLÉS À RETENIR}

\subsection*{Chronologie}
1971: Creeper | 1977: RSA | 1983: TCP/IP | 1987: Antivirus | 1988: Morris Worm + CERT

\subsection*{Triade CIA}
Confidentialité (accès) | Intégrité (modification) | Disponibilité (accessibilité)

\subsection*{Détection}
Signature (connues) | Comportementale (APT) | IA/ML (patterns)

NIDS (réseau) | HIDS (hôte) | IPS (bloque) | SIEM (corrélation)

\subsection*{Réponse NIST (6 phases)}
Préparation - Détection - Confinement - Éradication - Récupération - REX

\subsection*{Kill Chain (7 phases)}
Reconnaissance - Weaponization - Delivery - Exploitation - Installation - C2 - Actions

\subsection*{Métriques}
MTTD (Time To Detect) | MTTR (Time To Respond)

\subsection*{Teams}
Blue (défense) | Red (attaque) | Purple (collaboration)

\end{multicols}
\end{document}
